\chapter{Bevezetés}
\label{ch:intro}

Az elmúlt években, a digitális technológiák fejlődésének eredményeként, jelentősen átalakultak a nyelvtanulás módszerei és segédeszközei. A mesterséges intelligencia gyors térnyerésének köszönhetően ma már számos online platform, applikáció és eszköz áll az idegen nyelveket elsajátítani vágyók rendelkezésére. Ezek akár hagyományos, akár modernebb, egyedibb stratégiákkal segíthetik a szókincs, a nyelvtan vagy a szövegértés fejlesztését. Ugyanakkor ezek a rendszerek legtöbbször szeparálva vannak jelen, külön applikáció szükséges szövegek elemzéséhez, szavak vagy írásjelek memorizálásához, nyelvi statisztikák követéséhez és megjelenítéséhez. Ez gyakran eredményezheti a tanulási élmény romlását, hiszen a tanuló nem feltétlenül tud teljes képet kapni fejlődéséről és jelenlegi nyelvtudási szintjéről.

Az ilyen rendszerek mögött gyakran természetes nyelvfeldolgozó (\textit{Natural Language Processing, NLP}) megoldások állnak, amelyek lehetővé teszik a szövegek automatikus elemzését és értelmezését. A felhőalapú NLP szolgáltatások, például a Google Cloud NLP, kész megoldásokat kínálnak, azonban ezek gyakran korlátozott testreszabhatóságot biztosítanak, különösen egyedi nyelvi metrikák vagy tanulási szempontok kezelésében.

A dolgozat megírásának ötlete részben személyes tapasztalatból született. Az angol nyelvet már gimnazista koromban elsajátítottam, főleg iskolai keretek között, és azóta is minden nap használom. A japán nyelvvel körülbelül két éve kezdtem el komolyabban foglalkozni, és már a tanulás kezdetén nyilvánvalóvá vált, hogy az egyedi írásrendszer és a nyelvi szerkezet sajátosságai miatt egészen más megközelítést igényel. Eddigi tanulmányaim során számos digitális eszközt próbáltam ki, köztük sokat a mai napig használok, például a ma már a legtöbb nyelvtanuló által ismert Anki alkalmazást, egy a kifejezések és szavak memorizálását segítő eszközt, amely szakdolgozatom legfőbb inspirációja volt. Az Anki és hasonló alkalmazások hatékonysága ellenére idővel azt tapasztaltam, hogy ideálisabb lenne a több, különböző funkciót ellátó alkalmazás helyett ezek egy rendszerben történő egyesítése. 

Hasonló funkciókat kínáló alkalmazások közé tartozik például a Bunpro, a Quizlet vagy a Duolingo, amelyek előre összeállított tananyagokra építenek. Ezek előnye az egyszerű használat és a gyors tanulási kezdés lehetősége, ugyanakkor kevésbé támogatják az egyedi tartalmak feldolgozását és a személyre szabott tanulási folyamat kialakítását.

Célom egy olyan nyelvi feldolgozó és tudásszint követő rendszer megtervezése és implementálása volt, amely angol és japán szövegek automatikus elemzésére épül. A feldolgozás során a rendszer a bemeneti szöveget mondatokra bontja, majd azokat tovább elemzi szavakra és (japán esetében) kanji karakterekre. A kinyert nyelvi egységek egy adatbázisba kerülnek, ahol kiegészülnek különböző tulajdonságokkal (pl. olvasatok, vonásszám, szóosztály), valamint hozzájuk rendelhető nyelvi nehézségi szint (angol szavaknál CEFR, japán elemeknél JLPT), illetve gyakorisági adatokkal.

A rendszer megvalósítása során ezért nem felhőalapú megoldásokra, hanem \textit{open source} könyvtárakra támaszkodtam. A tokenizálás és nyelvi feldolgozás olyan eszközök segítségével történik, mint például a spaCy és a MeCab, amelyek lehetővé teszik a feldolgozási folyamatok teljes körű testreszabását, valamint az egyedi nyelvi attribútumok hozzárendelését.

Az elmentett adatokból a tanuló különböző szótárakat hozhat létre, amelyek az Anki alkalmazás által népszerűsített, ismétlésen alapuló tanulás elvét követve segítik a memorizálást. Ilyenkor egy napi korlát alapján új szavakat vagy írásjeleket tud elsajátítani, majd a rendszer a felhasználói visszajelzés alapján frissíti ezek tudásszintjét, és egy későbbi dátumra helyezi ismétlési célzattal. Mind a létrehozott szótárakban, mind az elemzett szövegekben mentett elemeket lehet exportálni Excel, CSV, vagy Quizlet által felismerhető formában, ezúton a rendszer más környezetekkel is összekapcsolható.

A LanDB előnye, hogy képes a szövegelemzés mellett folyamatos nyelvi támogatást biztosítani a felhasználó nyelvi képességeit és szükségleteit figyelembe véve. Nyelvtanulás során fontos, hogy a szavakat és kifejezéseket kontextuson belül lássuk, segítve ezzel a memorizálást és az adott nyelven való gondolkodásmód fejlődését.
